\section{Mechanisms and Heterogeneity}
\label{sec:mechanism}

\subsection{Valuation versus Profitability}

Table \ref{tab:main} shows a combination that is central to the interpretation: gross margin declines by 2.45 percentage points after the shock, while ROA and expense ratios do not change significantly. This pattern implies that the privacy shock compresses the profitability of revenue rather than the scale of operating costs. A natural explanation is that third-party data support pricing, targeting, and conversion. When the data become less available, the same revenue requires more acquisition effort or generates lower conversion, which lowers gross margin. Firms simultaneously cut low-efficiency advertising and administrative spending, which keeps ROA stable in the short run. The decline in data rents therefore appears in the quality of revenue rather than in the level of expenses.

The selling-expense ratio is consistent with this interpretation. If the shock operated mainly through higher customer-acquisition costs, the selling-expense ratio should rise. The estimated effect is $-0.0093$ (standard error 0.0077), negative and insignificant. Firms did not simply spend more on advertising; they adjusted the composition of spending, reduced low-efficiency channels, or moderated acquisition intensity. A falling gross margin without a rising selling-expense ratio points to a decline in data-driven pricing and matching, not to a pure cost shock.

\subsection{B2B versus B2C Firms}

Table \ref{tab:hetero} reports the heterogeneity results. B2B data-service firms—IT services, software, and communication services—experience a 14.2\% decline in log Tobin's Q (standard error 4.4\%), significant at the 1\% level. Consumer internet firms—advertising, gaming, e-commerce, digital media, film, publishing, broadcasting, and education—experience a 2.0\% decline (standard error 5.2\%), which is not significant. The difference is economically large: the point estimate for B2B firms is about seven times that for consumer firms, and a pooled interaction model yields a difference significant at the 10\% level ($p=0.107$). The top-third versus bottom-third comparison yields a 10.7\% decline in log Tobin's Q, a 0.329 decline in Tobin's Q, and a 3.41 percentage-point decline in gross margin.

This is the paper's central mechanism result. Third-party data dependence is not evenly distributed across digital firms. B2B data-service firms obtain data through external procurement, client authorization, cross-organizational sharing, and project-based delivery. PIPL's consent, minimization, and sharing rules directly raise the compliance cost of these activities, while ATT weakens their ability to provide cross-app attribution and targeted delivery to clients. Consumer internet firms own accounts, transaction records, first-party apps, and membership relationships; first-party data can substitute for part of the lost third-party data. Their valuations are therefore less sensitive to the privacy shock.

The result also clarifies what a ``data asset'' is. Data are not homogeneous bits. A unit of first-party data and a unit of third-party data have different substitutability under privacy regulation, and therefore different value. Treating all digital firms as a single data-asset category obscures the structural difference between B2B and B2C firms. The TDDI identifies this difference because it measures dependence on external data, not the level of digitalization.

\subsection{First-Party Data Capability}

Table \ref{tab:hetero} also reports the FPDCI split. High-FPDCI firms experience a 12.6\% decline in log Tobin's Q (standard error 5.5\%), while low-FPDCI firms experience a 9.5\% decline (standard error 5.3\%); both are significant, and the point estimate for the high-FPDCI group is slightly larger. The App-exposure split is similarly muted: high-app-exposure firms decline by 8.7\% (standard error 4.4\%), and low-app-exposure firms decline by 9.4\% (standard error 6.7\%). These results do not support a simple substitution story in which first-party data fully insulate firms from the shock.

A more plausible interpretation is that first-party data capability and third-party data dependence coexist. High-FPDCI firms are often data-intensive firms: they own user accounts and transaction data, but they also purchase external data to extend coverage, improve models, and conduct cross-platform advertising. When privacy regulation tightens, first-party data can buffer part of the shock, but it cannot replace the breadth and cross-domain reach of third-party data. The two data types are complements as much as substitutes. This finding qualifies the popular view that private-domain traffic is a safe harbor under privacy regulation: private data are valuable, but they do not make firms immune to the contraction of the third-party data ecosystem.

\begin{table}[htbp]
\centering
\small
\caption{Heterogeneity by firm type, app exposure, and first-party data capability.}
\label{tab:hetero}
\begin{threeparttable}
\begin{tabular}{lccc}
\toprule
 & Log Tobin's Q & Tobin's Q & Gross margin \\
\midrule
B2B firms & -0.1424*** (0.0438) & -0.4902*** (0.1637) & -0.0339*** (0.0121) \\
Consumer firms & -0.0197 (0.0523) & -0.1556 (0.1849) & -0.0095 (0.0198) \\
Low app exposure & -0.0942 (0.0671) & -0.2640 (0.1996) & -0.0155 (0.0192) \\
High app exposure & -0.0866* (0.0444) & -0.3972** (0.1714) & -0.0191 (0.0156) \\
Low FPDCI & -0.0952* (0.0525) & -0.4518** (0.1923) & 0.0006 (0.0179) \\
High FPDCI & -0.1261** (0.0553) & -0.3572* (0.2064) & -0.0505*** (0.0189) \\
\bottomrule
\end{tabular}
\begin{tablenotes}[flushleft]
\footnotesize
\item Notes: Each cell reports the coefficient on High TDDI $\times$ Post, with firm-clustered standard errors in parentheses. B2B firms include IT services, software development, and communication services. Consumer firms include advertising, gaming, e-commerce, digital media, film, publishing, broadcasting, and education. App exposure and FPDCI are split at the median. The top-third versus bottom-third comparison is reported in the text.
\end{tablenotes}
\end{threeparttable}
\end{table}


\subsection{Robustness}

Table \ref{tab:robust} reports 11 robustness checks. First, replacing the baseline TDDI with a narrow measure based only on category A (third-party data acquisition) yields a coefficient of $-0.0824$ (standard error 0.0355), significant at the 5\% level. Second, replacing the TDDI with the generic digitalization index yields $-0.0214$ (standard error 0.0193), which is insignificant. Third, in a horse race with both high TDDI and high generic digitalization, the TDDI coefficient is $-0.0748$ (standard error 0.0363) and the generic coefficient is $-0.0396$ (standard error 0.0349), with only the former significant. Fourth, excluding 2020 yields $-0.0932$. Fifth, restricting the sample to firms with at least 20 quarters yields $-0.0976$. Sixth, excluding firms in the top 1\% of assets yields $-0.0945$. Seventh, controlling for a 2020--2021 pandemic interaction yields $-0.1027$. Eighth, two-way clustering by firm and quarter yields a standard error of 0.0339 and leaves the significance unchanged. Ninth, winsorizing at the 5th and 95th percentiles yields $-0.0777$. Tenth, the top-third versus bottom-third comparison yields $-0.1067$. The wild cluster bootstrap $p$-value for the baseline coefficient is 0.007. All checks point in the same direction: the result is not driven by a single measure, subsample, extreme value, or clustering choice.

\begin{table}[htbp]
\centering
\small
\caption{Robustness checks for log Tobin's Q.}
\label{tab:robust}
\begin{threeparttable}
\begin{tabular}{lcccc}
\toprule
 & Coefficient & Standard error & $p$-value & Observations \\
\midrule
Baseline: high TDDI & -0.0917*** & (0.0343) & 0.007 & 6,900 \\
Narrow TDDI measure & -0.0824** & (0.0355) & 0.020 & 6,900 \\
Placebo: generic digitalization & -0.0214 & (0.0193) & 0.268 & 6,900 \\
Horse race: TDDI and generic digitalization & -0.0748** & (0.0363) & 0.039 & 6,900 \\
Excluding 2020 & -0.0932*** & (0.0339) & 0.006 & 5,733 \\
Balanced panel (at least 20 quarters) & -0.0976*** & (0.0344) & 0.005 & 6,458 \\
Excluding the top 1\% by assets & -0.0945*** & (0.0351) & 0.007 & 6,831 \\
Controlling for the 2020--2021 pandemic interaction & -0.1027*** & (0.0396) & 0.009 & 6,900 \\
Two-way clustering (firm and quarter) & -0.0917*** & (0.0339) & 0.007 & 6,900 \\
Winsorizing at the 5th and 95th percentiles & -0.0777** & (0.0329) & 0.018 & 6,900 \\
Top third versus bottom third & -0.1067** & (0.0443) & 0.016 & 4,631 \\
\bottomrule
\end{tabular}
\begin{tablenotes}[flushleft]
\footnotesize
\item Notes: All columns use log Tobin's Q as the dependent variable and include firm fixed effects, industry-by-quarter fixed effects, and controls. The wild cluster bootstrap $p$-value for the baseline coefficient is 0.007 (999 replications). The generic digitalization index is the standardized frequency of category F terms; it should be insignificant if the TDDI captures a distinct construct.
\end{tablenotes}
\end{threeparttable}
\end{table}

